\documentclass[fontset=ubuntu]{ctexart}
\usepackage[showframe,margin=1in]{geometry}
\usepackage{amsmath,amsfonts}
\usepackage{enumitem}
\usepackage{zhlipsum}
\newcounter{question}%
\NewDocumentEnvironment{question}{ O +b }{%
    \refstepcounter{question}%
    \noindent%
    \begin{minipage}[t]{1.5cm}%
    \renewcommand{\arraystretch}{1.25}%
    \begin{tabular}{|p{1cm}|}%
        \hline%
        得\enspace 分\\%
        \hline\\\hline%
    \end{tabular}%
    \end{minipage}%
    \begin{minipage}[t]{\dimexpr\textwidth-1.5cm\relax}%
    \setlength{\parindent}{0pt}%
    \arabic{question}、%
    \IfNoValueF{#1}{（#1分）}#2%
    \end{minipage}%
}%
{\bigskip}%
\begin{document}
\zhlipsum[2]

\noindent
\begin{minipage}[t]{1.5cm}
\renewcommand{\arraystretch}{1.25}
\begin{tabular}{|p{1cm}|}
    \hline
    得\enspace 分\\
    \hline\\\hline
\end{tabular}%
\end{minipage}%
\begin{minipage}[t]{\dimexpr\textwidth-1.5cm\relax}
\setlength{\parindent}{0pt}
3、（14分）设$X_1,\ldots,X_n$为来自均匀分布$U(\theta_1,\theta_2)$的i.i.d.样本, 其中$\theta_1<\theta_2$为未知参数.
\begin{enumerate}[label=(\roman*)]
    \item 求$(\theta_1+\theta_2)/2$的矩估计.
    \item 如果$T(X)=(X_{(1)},X_{(n)})$是完备的, 则证明$(X_{(1)}+X_{(n)})/2$是$(\theta_1+\theta_2)/2$的UMVUE.
\end{enumerate}%
\end{minipage}%
\bigskip

\zhlipsum[2]
\setcounter{question}{3}
\begin{question}[14]%
    设$X_1,\ldots,X_n$为来自均匀分布$U(\theta_1,\theta_2)$的i.i.d.样本, 其中$\theta_1<\theta_2$为未知参数.
    \begin{enumerate}[label=(\roman*)]
        \item 求$(\theta_1+\theta_2)/2$的矩估计.
        \item 如果$T(X)=(X_{(1)},X_{(n)})$是完备的, 则证明$(X_{(1)}+X_{(n)})/2$是$(\theta_1+\theta_2)/2$的UMVUE.
    \end{enumerate}
\end{question}

\zhlipsum[2]
\end{document}